% Multi-panel hybrid figure skeleton — raster slots + pure-TikZ panels.
% Adapted from the proven atom-ion figure-1 prototype.
%
% Workflow:
%   1. Generate raster components with scripts/hybrid_gen.py into assets/
%   2. Adjust \includegraphics paths + panel labels below
%   3. Tune coordinates via preview iterations
\documentclass[border=4pt]{standalone}

\usepackage{tikz}
\usepackage{graphicx}
\usepackage{amsmath,amssymb}
\usepackage{siunitx}
\usepackage{braket}
\usepackage{xcolor}

\usetikzlibrary{positioning,arrows.meta,calc,decorations.pathmorphing,
                decorations.markings,backgrounds}

\definecolor{oiBlue}{HTML}{0072B2}
\definecolor{oiVerm}{HTML}{D55E00}
\definecolor{oiOrange}{HTML}{E69F00}
\definecolor{oiPurple}{HTML}{8E44AD}
\definecolor{oiGreen}{HTML}{009E73}
\definecolor{labelGray}{HTML}{444444}

\colorlet{atomC}{oiBlue}
\colorlet{ionC}{oiVerm}
\colorlet{laserC}{oiOrange}
\colorlet{accentC}{oiPurple}
\colorlet{highlightC}{oiGreen}

\begin{document}
\begin{tikzpicture}[font=\sffamily]

  %% ====================================================
  %% Panel (a) — raster + TikZ labels
  %% ====================================================
  \begin{scope}[local bounding box=panelA]
    % Raster slot 1
    % \node[inner sep=0pt] (obj1) at (0,0) {\includegraphics[height=4cm]{assets/obj1.png}};

    % TikZ labels / arrows overlaid
    % \node[atomC, font=\small\bfseries] at ($(obj1)+(0,-1.5)$) {label};

    \node[font=\large\bfseries] at (-2.5, 2.3) {(a)};
  \end{scope}

  %% ====================================================
  %% Panel (b) — pure TikZ
  %% ====================================================
  \begin{scope}[xshift=9cm, local bounding box=panelB]
    % Draw energy levels, arrows, kets here
    \node[font=\large\bfseries] at (-0.5, 2.3) {(b)};
  \end{scope}

  %% ====================================================
  %% Panel (c) — pure TikZ
  %% ====================================================
  \begin{scope}[xshift=14cm, local bounding box=panelC]
    % Force diagram / phase-space / schematic conclusion
    \node[font=\large\bfseries] at (-0.5, 2.3) {(c)};
  \end{scope}

\end{tikzpicture}
\end{document}
